% Autor: Marek Vandík, 13. 4. 2026

\documentclass[a4paper,11pt]{article}

% preambule

\usepackage[czech]{babel}
\usepackage[left=2cm,text={17cm,24cm},top=3cm]{geometry}
\usepackage[hidelinks]{hyperref}
\usepackage{times}
\usepackage{multirow}
%\usepackage[czech,linesnumbered,longend,noline,ruled]{algorithm2e}
\usepackage{graphics}
\usepackage{picture}
\usepackage{pdflscape}

% bibtex + babel - nazev citaci
\addto\captionsczech{\renewcommand{\refname}{Literatura}}

% titulni strana

\begin{document}

\begin{titlepage}
\begin{center}
	\Huge{\textsc{Vysoké učení technické v~Brně}} \\
	\huge{\textsc{Fakulta informačních technologií}} \\
	\vspace{\stretch{0.382}}
		\LARGE{Přenos dat, počítačové sítě a~protokoly\,--\,projekt} \\
		\Huge{Analýza transportních protokolů a~algoritmů pro~řízení zahlcení} \\
	\vspace{\stretch{0.618}}
	\Large{\today \hfill Marek Vandík}
\end{center}
\end{titlepage}

\tableofcontents
\newpage

% TODO navrhované optimalizace
\section{Úvod}
Projekt se zabývá analýzou chování algoritmů řízení zahlcení při~různých síťových podmínkách pro~transportní
protokol~TCP a~dále tyto~přístupy porovnává s~protokolem~QUIC.

Pro~účely projektu bylo vytvořeno simulační prostředí na~bázi
kontejnerizace, vizte podkapitolu~\ref{subsec-containers}. % TODO github
V~kapitole~\ref{sec-transport} naleznete vysvětlení, jak~jsou dané protokoly a~jejich přístupy k~řízení zahlcení specifikovány.
Kapitola~\ref{sec-experiment} uvádí, jaká topologie sítě byla zvolena, jak~experiment probíhal včetně popisu zvolené metodiky, a~s~jakými sítěmi se pracovalo.
Výsledky jsou uvedeny v kapitole~\ref{sec-result} včetně vizualizací klouzavého okna odesilatele, odezvy, dosažené rychlosti,
a~porovnání algoritmů pro~dané případy použití. Ve~stejné kapitole je~taktéž porovnán protokol~TCP s~protokolem~QUIC.

\section{Zkoumané transportní protokoly} \label{sec-transport}
\subsection{Transport s TCP}
\subsection{Transport s QUIC}
\subsection{Algoritmy řízení zahlcení}
\subsubsection{Reno}
Bleh blah \cite{rfc5681}.
\subsubsection{Cubic}
Blih bleh \cite{rfc9438}.
\subsubsection{BBR}

\section{Experiment} \label{sec-experiment}
\subsection{Prostředí simulace} \label{subsec-containers}
\subsection{Topologie sítě a její parametry}
\subsubsection{Bezdrátová síť s velkou ztrátovostí (Wi-Fi)}
\subsubsection{Satelitní síť s vyšší odezvou (Starlink)}
\subsubsection{Síť s příliš velkou frontou na směrovači (Bufferbloat)}
\subsection{Metodika měření}

\section{Výsledky a zhodnocení} \label{sec-result}
\subsection{Vizualizace zahlcení}
\subsubsection{Bezdrátová síť s velkou ztrátovostí (Wi-Fi)}
\subsubsection{Satelitní síť s vyšší odezvou (Starlink)}
\subsubsection{Síť s příliš velkou frontou na směrovači (Bufferbloat)}
\subsection{Porovnání algoritmů řízení zahlcení}
\subsection{TCP vs. QUIC}

\section{Závěr}
Blih bleh blah

% bibtex - styl Pysny/czplain z sablony zaverecnych praci
\clearpage
\phantomsection
\addcontentsline{toc}{section}{Literatura}
\begin{flushleft}
	\bibliographystyle{czplain}
	\bibliography{citace}
\end{flushleft}

\end{document}
